\section{Raw balance and stability without a variance floor}
\subsection*{A finite moment criterion}
Let $E$ be a finite-dimensional real ordered space with generating positive cone and a strictly positive mass functional $u$. Its normalized positive base $S=\{s\ge0:u(s)=1\}$ is compact. This is a verification device for the following proposition; Theorem~\ref{thm:main} does not assume a finite-dimensional physical space.

For each of finitely many allowed within-call commands $j$, suppose the raw joint report--audit law depends linearly and positively on $s\in S$. After the declared report transformation, and a task isometry $T_j:H\to H$ with $T_jH=H$, write $\widetilde Z=T_j^{-1}Z$. The following vector measures are linear functionals of $s$:
\[
 A_j(D)(s)=\E_s[\one_{\{x\in D\}}],\quad
 B_j(D)(s)=\E_s[\one_{\{x\in D\}}\widetilde Z],\quad
 C_j(D)(s)=\E_s[\one_{\{x\in D\}}\|\widetilde Z\|^2].
\]
Choose an origin in $\operatorname{aff}(H)$ and orthonormal coordinates on its $k_0=\dim\operatorname{aff}(H)$-dimensional direction space. Here $B_j$ has $k_0$ coordinates; the written squared norms are in these translated coordinates. Translation leaves all compared squared losses unchanged. Fix a genuinely executed reference law $K_0$ and its moments \eqref{eq:derived}.

\begin{proposition}[Raw task-moment balance criterion]\label{prop:moment}
A fixed distribution $\lambda\in\Delta(J)$ balances report and squared-loss measures for every normalized input $s$ if and only if, as $E^*$-valued measures,
\begin{align}
 \sum_j\lambda_j A_j(dx)&=\nu(dx)u,\label{eq:moment0}\\
 \sum_j\lambda_j B_j(dx)&=f(x)\nu(dx)u,\label{eq:moment1}\\
 \sum_j\lambda_j C_j(dx)&=(w(x)+\|f(x)\|^2)\nu(dx)u.\label{eq:moment2}
\end{align}
Equality on a countable determining algebra of report events suffices. The set of such $\lambda$ is compact and convex. If it is empty, finitely many scalar moment constraints already separate all mixtures in $\Delta(J)$ from the required identities.
\end{proposition}
\begin{proof}
Expand $\|\widetilde Z-r\|^2=\|\widetilde Z\|^2-2r\cdot\widetilde Z+\|r\|^2$. The three identities give \eqref{eq:mass-balance}--\eqref{eq:loss-balance}, using that $T_j$ is an isometry. Conversely, the mass identity and equality of this quadratic polynomial for $r$ in the affine hull of $H$ determine the first moments in that affine hull and the squared-norm moment. One may first translate the affine hull to the origin and use an orthonormal coordinate system on it; orthogonal coordinates of $Z$ are then fixed and carry no additional constraint. Equality for all $s\in S$ is equality as linear functionals because the positive cone is generating. Uniqueness of finite signed measures extends equality from a determining algebra to all Borel sets. Each coordinate constraint is affine and closed in the finite simplex. Compactness and the finite-intersection property give a finite inconsistent subfamily if no mixture exists. Its affine image of the simplex is compact and convex and omits the prescribed moment vector. Finite-dimensional strict separation proves the last assertion. This is a certificate of impossibility for the declared command family, not for arbitrary causal algorithms.
\end{proof}

The criterion is insensitive to density zeros or changes in the rank of a posterior. It also shows why report-only averaging is insufficient. A report identically equal to zero has the same law for hidden states $e=\pm1$, but an audit $Z=e$ has loss $(e-r)^2$. Without balancing its first moment there is no common $d_r$ valid for all outputs, despite perfect report balance.

A finite physical randomization satisfying $\sum_j\lambda_jG_j(s)=\omega u(s)$ is one sufficient source of \eqref{eq:moment0}--\eqref{eq:moment2} for any subsequent fixed instrument. More generally, it suffices to annihilate input dependence only on the report, first-moment and squared-norm functionals; the entire hidden state need not be erased. These are linear moment identities, not an assertion that a Fisher matrix, covariance or fixed-rank parametrization is primitive. Randomization itself is standard \cite{Luschgy}, \cite{AMTW}; the stopped conclusion follows only after the actual witness, paid acquisition and exported readout conditions have also been verified.

\subsection*{Integrated predictive-update stability}
The useful continuity is that of unnormalized instruments. Uniform pointwise continuity of a normalized posterior can fail at zero-mass reports.

\begin{proposition}[Integrated stability through singular reports]\label{prop:instrument}
Let a classical stochastic instrument, or a finite-dimensional completely positive trace-preserving instrument, map two normalized inputs $s,t$ to subnormalized post-report states $A_x,B_x$ with respect to a common finite dominating measure for these two inputs. This measure may be chosen pairwise; the inequalities are unchanged by replacing it with any equivalent dominating representation. Put $p_x=u(A_x)$, $q_x=u(B_x)$, and normalize on positive-mass reports. In the classical $\ell^1$ norm or quantum trace norm,
\begin{align}
 \int\|A_x-B_x\|_1\,dx&\le\|s-t\|_1,\label{eq:unnorm-stable}\\
 \|p-q\|_{\TV}&\le\tfrac12\|s-t\|_1,\label{eq:report-stable}\\
 \int\min(p_x,q_x)\left\|\frac{A_x}{p_x}-\frac{B_x}{q_x}\right\|_1dx
 &\le2\|s-t\|_1.\label{eq:weighted-stable}
\end{align}
Zero-mass terms are defined to contribute zero. The same inequalities apply to a complete finite causal continuation instrument under the same fixed parameter-independent policy. If its stopped scored instrument terminates almost surely under both inputs, the unnormalized contraction and report bound extend to that stopped instrument.
\end{proposition}
\begin{proof}
In the quantum case write the Hermitian difference $s-t=D_+-D_-$ with positive parts. Positivity, the triangle inequality and preservation of integrated trace give
$\int\|\mathcal I_x(D_+-D_-)\|_1dx\le\tr D_++\tr D_-=\|s-t\|_1$.
The same Jordan-decomposition argument works for classical signed measures. Taking masses proves \eqref{eq:report-stable}. When $p\le q$,
\[
 p\|A/p-B/q\|_1=\|A-(p/q)B\|_1
 \le\|A-B\|_1+(q-p).
\]
Interchange $A,B$ for the other case and integrate, using $|p-q|\le\|A-B\|_1$. Composition of legal physical instruments remains positive and trace preserving, so the finite-continuation assertion follows. For stopped paths, sum over commitment time, first to a finite horizon while retaining the surviving output. This is a positive mass-preserving instrument. Almost-sure termination under both inputs lets the surviving masses tend to zero; monotone summation gives the stopped assertion. No bound on expected stopping time follows from this argument; it requires a separate cost certificate.
\end{proof}

This stability controls the \emph{actual report mass} together with the posterior, including pure states and likelihood zeros. It is not a contraction factor strictly below one, and no normalized posterior is stored by the finite-state upper. In particular a pointwise variance-floor constant diverging along nearly pure histories does not force the risk constant of Theorem~\ref{thm:main} to diverge.
